\section{Introduction}

Accurate identification and integrated management of crop pests and diseases constitute a cornerstone of sustainable agriculture, directly influencing smallholder farmer income, regional food security, crop yield stability, and the environmental footprint of agrochemical use \cite{jearanaiwongkul2021riceman}. In recent years, artificial intelligence has been increasingly deployed to assist agricultural practitioners in diagnostic and advisory tasks. Contemporary technological approaches broadly bifurcate into two paradigms: cloud-based large language models (LLMs), exemplified by industrial systems such as the Shennong Plant-Protection multi-modal model \cite{wang2024shennong} and the XiongXiaoNong agricultural LLM ecosystem, and edge-oriented lightweight inference pipelines, typically based on convolutional neural networks (CNNs) or vision transformers deployed on embedded devices \cite{chatdemeter2025}. While both directions have yielded promising accuracy metrics on benchmark datasets, they share a deeper architectural limitation that hinders scalable, equitable, and cost-effective deployment across the world's diverse agro-ecological contexts.

A fundamental pain point common to both paradigms is the tight embedding of domain logic within model weights, inference code, or hard-coded ontologies \cite{alharbi2024agrodss}. When a system originally designed for rice blast detection needs to be repurposed for wheat rust, citrus canker, or tomato early blight, developers are often compelled to rewrite semantic web rules, retrain vision backbones, redesign agent orchestration chains, or reconstruct knowledge graphs from scratch \cite{agenticgraphrag2025}. This high migration cost renders existing solutions essentially ``vertical applications'' rather than horizontally extensible platforms. Moreover, cloud-centric architectures introduce substantial latency, recurring subscription costs, privacy risks surrounding farm-level data, and availability concerns in rural regions where network connectivity is intermittent, expensive, or entirely absent during critical growing seasons \cite{kleppmann2021localfirst}. Edge inference boxes, by contrast, typically offer only single-function model execution---such as image classification---without the higher-level task planning, structured knowledge retrieval, tool coordination, and explainable reasoning capabilities required of a fully fledged intelligent agent.

These empirical observations lead to three specific gaps in the current literature and practice. First, there is a paucity of general-purpose agent frameworks that can adapt to disparate agricultural sub-domains---ranging from row crops to orchard systems and from greenhouse vegetables to open-field cereals---through declarative configuration alone, without any modification to the underlying source code. Second, although the Model Context Protocol (MCP) has recently emerged as a promising open standard for augmenting LLMs with external tools \cite{anthropic2024mcp}, its adaptation to agriculture---including tool-registration schemas for pest ontologies, context-exchange formats for farm-level knowledge, and offline transport mechanisms---remains largely unexplored in published research. Third, truly local-first, offline-capable agricultural knowledge-management agents that integrate task planning, knowledge-graph retrieval, and tool-use reasoning within a unified edge-native runtime have yet to be demonstrated in the agricultural informatics literature.

To address these gaps, the present study establishes three concrete research objectives (ROs):
\begin{enumerate}[label=RO\arabic*:,leftmargin=*]
    \item To verify the technical feasibility of coupling a Rust-based agent framework (Clarity) with a local knowledge-base backend (devbase) via the MCP protocol over stdio transport, thereby decoupling the application layer from the domain knowledge layer.
    \item To propose and empirically validate a ``Domain-as-Config'' methodology that enables zero-code adaptation of a generic agent kernel to agricultural knowledge-management scenarios through external TOML configuration files.
    \item To evaluate the performance, reliability, and end-user usability of the resulting local-first architecture in offline agricultural knowledge-retrieval and pest-diagnosis tasks.
\end{enumerate}

Against this backdrop, the principal contributions of this work are threefold. First, we present the first MCP-based, configuration-driven agent architecture specifically oriented toward agricultural knowledge management, in which crop-specific personas, tool inventories, and knowledge-base endpoints are injected at runtime rather than hard-coded at build time. Second, we implement a complete, standards-compliant Rust MCP client/server communication mechanism, including Content-Length-framed JSON-RPC~2.0 message exchange, dynamic tool-schema discovery, and robust stdio transport alignment, all rigorously verified through integration testing against the official MCP specification. Third, through systematic integration testing and performance benchmarking, we demonstrate a viable migration path from a generic software-engineering assistant to an agricultural domain expert purely via declarative configuration files, thereby reducing cross-crop migration time from hours of coding to minutes of editing while maintaining sub-second query latency on local SQLite backends.

The remainder of this paper is organized as follows. Section~\ref{sec:related-work} reviews related work in agricultural expert systems, LLMs and knowledge graphs, MCP protocols, and Rust agent frameworks. Section~\ref{sec:methodology} details the proposed architecture, the MCP adaptation mechanism, and the Domain-as-Config paradigm. Section~\ref{sec:experiments} presents the experimental design, baseline comparisons, and results. Section~\ref{sec:discussion} discusses the implications, limitations, and future research directions. Finally, Section~\ref{sec:conclusion} summarizes the principal findings and discusses their broader significance for the future of agricultural software architecture.
